

\hypertarget{fixture-mla9-parity}{%
\section{Fixture: mla9 parity}\label{fixture-mla9-parity}}

The foundational treatment of pragmatic markers appears in Smith, whose
framework has shaped later work (Smith 22). Classroom-discourse data in
Jones and Cruz produced complementary findings (Jones and Cruz 84--86).
The urban multilingual synthesis broadened the empirical base (Able et
al.).

Smith argues that such markers serve multiple pragmatic functions. A
consolidated framing across the three strands is now available (Smith;
Jones and Cruz).

Institutional guidance on multi-author citation resolves common
ambiguities (Modern Language Association). Undated industry commentary
touches on adjacent practice (Perennial).

\hypertarget{bibliography}{%
\section*{Works Cited}\label{bibliography}}
\addcontentsline{toc}{section}{Works Cited}

\hypertarget{refs}{}
\begin{CSLReferences}{1}{0}
\leavevmode\vadjust pre{\hypertarget{ref-able-baker-chen-2019}{}}%
Able, Amir, et al. {``Multilingual Practices in Urban Contexts.''}
\emph{Language in Society}, vol. 48, no. 1, 2019, pp. 55--80.

\leavevmode\vadjust pre{\hypertarget{ref-jones-cruz-2015}{}}%
Jones, Robert, and Elena Cruz. {``Code-Switching in Classroom
Discourse.''} \emph{Applied Linguistics}, vol. 36, no. 2, 2015, pp.
84--102.

\leavevmode\vadjust pre{\hypertarget{ref-mla-2021}{}}%
Modern Language Association. {``How Do I Cite a Source with Multiple
Authors?''} \emph{MLA Style Center}, 2021,
\url{https://style.mla.org/citing-multiple-authors/}.

\leavevmode\vadjust pre{\hypertarget{ref-perennial-nd}{}}%
Perennial, Paul. {``Undated Commentary.''} \emph{Perennial Blog},
\url{https://example.com/undated}.

\leavevmode\vadjust pre{\hypertarget{ref-smith-2010}{}}%
Smith, Jane A. {``Pragmatic Markers in Contact Situations.''}
\emph{Journal of Pragmatics}, vol. 42, no. 3, 2010, pp. 12--34,
\url{https://doi.org/10.1016/j.pragma.2010.01.005}.

\end{CSLReferences}

